\documentclass[a4paper,12pt]{article}
\usepackage[T2A]{fontenc}
\usepackage[utf8]{inputenc}
\usepackage[russian]{babel}
\usepackage{amsmath,amssymb,mathtools}
\usepackage{geometry}
\geometry{left=23mm,right=21mm,top=22mm,bottom=23mm}
\usepackage{microtype}
\usepackage{icomma}
\usepackage{booktabs,tabularx,longtable}
\usepackage{enumitem}
\usepackage{needspace}
\usepackage{tikz}
\usetikzlibrary{arrows.meta,positioning,shapes.geometric,calc}
\usepackage{xcolor}
\usepackage{fancyhdr}
\usepackage{setspace}
\usepackage{hyperref}
\definecolor{accent}{HTML}{008888}
\definecolor{darkaccent}{HTML}{005D64}
\definecolor{ink}{HTML}{252B30}
\definecolor{muted}{HTML}{65737A}
\hypersetup{colorlinks=true,linkcolor=darkaccent,urlcolor=darkaccent}
\setstretch{1.22}
\setlength{\parindent}{0pt}
\setlength{\parskip}{5pt}
\setlist[itemize]{leftmargin=7mm,itemsep=3pt,topsep=3pt}
\setlist[enumerate]{leftmargin=8mm,itemsep=6pt,topsep=5pt}
\renewcommand{\arraystretch}{1.24}
\pagestyle{fancy}
\fancyhf{}
\fancyhead[L]{\small\textcolor{muted}{Оптимизация и производная}}
\fancyhead[R]{\small\textcolor{muted}{Матвей Горбачев}}
\fancyfoot[C]{\small\textcolor{muted}{\thepage}}
\renewcommand{\headrulewidth}{0.3pt}
\renewcommand{\footrulewidth}{0pt}
\setlength{\headheight}{15pt}
\newcommand{\key}[1]{\textcolor{darkaccent}{\textbf{#1}}}
\newcommand{\checkit}[1]{\par\noindent$\square$\enspace #1\par}
\newcommand{\ex}[1]{\textbf{Пример.}~#1}
\newcommand{\task}[2]{\Needspace{12\baselineskip}\par\medskip\noindent{\color{darkaccent}\bfseries Задача~#1.}~\textbf{#2}\par\smallskip}
\begin{document}
\thispagestyle{empty}
\begin{center}
\vspace*{27mm}
{\large\textcolor{darkaccent}{МАТЕМАТИКА}}\par
\vspace{15mm}
{\huge Чек-лист}\par
\vspace{14mm}
{\LARGE\textcolor{accent}{Оптимизация сложных функций}}\par
\vspace{3mm}
{\Large\textcolor{accent}{Производная, преобразования и проверка}}\par
\vspace{20mm}
{\normalsize Подготовка к ЕГЭ по профильной математике}\par
\vfill
{\large Матвей Горбачев}\par
\vspace{12mm}
{\normalsize 12 ноября 2026 г.}\par
\vspace{13mm}
{\small Лёвин Артём Александрович}\par
{\small эксклюзивно для Матвея Горбачева}\par
\vspace{19mm}
\begin{tikzpicture}
\draw[accent,line width=.9pt] (-3.3,0) .. controls (-2.25,.47) and (-1.55,-.40) .. (0,0) .. controls (1.30,.42) and (2.25,-.44) .. (3.3,0);
\end{tikzpicture}
\vspace{4mm}
\end{center}
\clearpage

\section*{1. Что требуется уметь}
В задачах на оптимизацию необходимо определить, при каком допустимом значении переменной функция достигает наибольшего или наименьшего значения. На занятии рассматривались сложные формулы из задач о вынужденных колебаниях и электрической цепи. Главный приём обеих задач: \key{упростить исходную дробь до дифференцирования}.

\checkit{Отличаю переменную от постоянных параметров. Если функция записана как $A(\tau)$, переменная~--- $\tau$; остальные величины постоянны, пока в условии не сказано иное.}
\checkit{Записываю допустимую область: время и циклическая частота положительны; знаменатель не обращается в нуль.}
\checkit{Умею делить числитель и знаменатель на одну и ту же положительную степень переменной.}
\checkit{Если функция имеет вид $K/\sqrt{D(x)}$ при $K>0$ и $D(x)>0$, ищу \textbf{минимум $D(x)$}, чтобы получить \textbf{максимум исходной функции}.}
\checkit{Нахожу производную, исследую её знак, проверяю границы области и физический смысл ответа.}
\checkit{Проверяю раскрытие скобок, знаки, степени в знаменателях и значения параметров.}

\section*{2. Главный принцип: заменить трудную функцию простой}
Пусть
\[
M(x)=\frac{K}{\sqrt{D(x)}},\qquad K>0,\quad D(x)>0.
\]
При увеличении $D(x)$ знаменатель возрастает, а значение дроби уменьшается. Поэтому
\[
\boxed{\quad M(x)\text{ наибольшая}\quad\Longleftrightarrow\quad D(x)\text{ наименьшая}.\quad}
\]
Это правило позволяет исследовать $D(x)$ вместо громоздкой функции с дробью и квадратным корнем. \textbf{Важно:} минимум должен \textbf{достигаться} на допустимой области; одно лишь стремление к некоторому пределу ещё не гарантирует существования максимума.

\ex{Для $M(x)=12/\sqrt{4+(x-3)^2}$ при $x>0$ подкоренное выражение наименьшее при $x=3$, потому что квадрат $(x-3)^2\ge 0$ и равен нулю ровно при $x=3$. Следовательно, максимальное значение $M$ равно $12/2=6$.}

\section*{3. Преобразования, которые сокращают вычисления}
Если $x>0$, то
\[
\frac{Kx^2}{\sqrt{B(x)}}
=\frac{K}{\sqrt{B(x)/x^4}}.
\]
Здесь знаменатель делится на $x^2$, поэтому под корнем появляется деление на $x^4$. Переход основан на равенстве $\sqrt{x^4}=x^2$.

\ex{Преобразуем функцию при $x>0$:}
\[
\frac{7x^2}{\sqrt{(16-x^2)^2+3x^2}}
=\frac{7}{\sqrt{\dfrac{(16-x^2)^2+3x^2}{x^4}}}
=\frac{7}{\sqrt{\dfrac{256}{x^4}-\dfrac{29}{x^2}+1}}.
\]
Проверка коэффициента: $(16-x^2)^2=256-32x^2+x^4$, а после прибавления $3x^2$ коэффициент при $x^2$ равен $-29$.

\key{Полезные правила}
\[
(a-b)^2=a^2-2ab+b^2,\qquad
\frac{x^p}{x^q}=x^{p-q}\quad(x\ne 0).
\]
\[
\left(\frac{a}{b}\right)'=\frac{a'b-ab'}{b^2}\quad(b\ne0),\qquad
(x^{-m})'=-m x^{-m-1}.
\]
Если знаменатель производной строго положителен на области, знак всей производной определяется только знаком числителя. Но положительность нужно \textbf{доказать} или получить из условия.

\clearpage
\section*{4. Колебания лайнера: разбор полного метода}
Амплитуда вынужденных колебаний задаётся формулой
\[
A(\tau)=\frac{k\theta_0\tau^2}
{\sqrt{(T^2-\tau^2)^2+n\tau^2T^2}},
\qquad \tau>0,\ T>0,\ k\theta_0>0,\ 0<n<0{,}5.
\]
Здесь $\tau$~--- период приходящей волны, $T$~--- период собственных колебаний, $n$~--- заданный параметр модели. Положительный множитель $k\theta_0$ остаётся постоянным при оптимизации по $\tau$.

\textbf{Шаг 1. Убираем переменную из числителя.} Делим числитель и знаменатель на $\tau^2$:
\[
A(\tau)=\frac{k\theta_0}{\sqrt{D(\tau)}},\qquad
D(\tau)=\frac{(T^2-\tau^2)^2+n\tau^2T^2}{\tau^4}.
\]
\textbf{Шаг 2. Аккуратно раскрываем квадрат и приводим подобные члены.}
\[
\begin{aligned}
D(\tau)&=\frac{T^4-2T^2\tau^2+\tau^4+nT^2\tau^2}{\tau^4}\\
&=\frac{T^4}{\tau^4}+\frac{(n-2)T^2}{\tau^2}+1.
\end{aligned}
\]
Коэффициент \textbf{$n-2$} отрицателен при $0<n<0{,}5$. Это существенно для дальнейшего исследования.

\textbf{Шаг 3. Находим производную удобным способом: по слагаемым.}
\[
\begin{aligned}
D'(\tau)&=-\frac{4T^4}{\tau^5}-\frac{2(n-2)T^2}{\tau^3}\\
&=\frac{2T^2\bigl((2-n)\tau^2-2T^2\bigr)}{\tau^5}.
\end{aligned}
\]
Знаменатель $\tau^5$ и множитель $2T^2$ строго положительны. Поэтому знак $D'$ задаёт выражение $(2-n)\tau^2-2T^2$.

\textbf{Шаг 4. Приравниваем производную к нулю и учитываем область.}
\[
(2-n)\tau^2-2T^2=0
\quad\Longrightarrow\quad
\tau^2=\frac{2T^2}{2-n}
\quad\Longrightarrow\quad
\boxed{\tau_*=T\sqrt{\frac{2}{2-n}}}.
\]
Отрицательное значение квадратного корня исключаем: период по смыслу положителен.

\textbf{Шаг 5. Проверяем экстремум по знакам.} При $0<\tau<\tau_*$ производная $D'$ отрицательна; при $\tau>\tau_*$ положительна. Значит, $D$ убывает, достигает минимума и возрастает. Амплитуда $A$, напротив, сначала возрастает, достигает максимума при $\tau_*$, затем убывает.

\textbf{Шаг 6. При необходимости находим наибольшую амплитуду.}
\[
D(\tau_*)=n-\frac{n^2}{4}
=n\left(1-\frac n4\right),\qquad
\boxed{A_{\max}=\frac{k\theta_0}{\sqrt{n(1-n/4)}}}.
\]
\key{Важное замечание о параметре:} если известны только $T=14$~с и $0<n<0{,}5$, то однозначного числового периода получить нельзя. Ответ зависит от $n$:
\[
\tau_*(n)=14\sqrt{\frac{2}{2-n}}\ \text{с},\qquad
14<\tau_*<\frac{28}{\sqrt3}\ \text{с}.
\]
Значения $n=0$ и $n=0{,}5$ \textbf{не входят} в открытый промежуток; подставлять $0{,}5$ как заданное значение без дополнительного условия нельзя.

\section*{5. Колебательный контур: напряжение на катушке}
Для последовательной электрической цепи дана амплитуда напряжения
\[
U_L(\omega)=\frac{UL\omega}
{\sqrt{R^2+\left(\omega L-\dfrac{1}{\omega C}\right)^2}},
\quad \omega>0,\quad U,L,C,R>0.
\]
Искомая переменная~--- циклическая частота $\omega$. Величины $U$, $L$, $C$ и $R$ считаются постоянными.

\textbf{Шаг 1. Делим числитель и знаменатель на $\omega>0$.}
\[
U_L(\omega)=\frac{UL}{\sqrt{D(\omega)}},\qquad
D(\omega)=\frac{R^2}{\omega^2}+
\left(L-\frac{1}{\omega^2 C}\right)^2.
\]
Теперь достаточно искать минимум $D(\omega)$.

\textbf{Шаг 2. Раскрываем квадрат разности.}
\[
\begin{aligned}
D(\omega)&=\frac{R^2}{\omega^2}
+L^2-\frac{2L}{\omega^2 C}+\frac{1}{\omega^4 C^2}\\
&=L^2+\frac{R^2-2L/C}{\omega^2}
+\frac{1}{\omega^4 C^2}.
\end{aligned}
\]
Проверяйте особенно тщательно: в слагаемых $2L/(\omega^2C)$ и $1/(\omega^4C^2)$ степени $\omega$ остаются \textbf{в знаменателе}.

\textbf{Шаг 3. Дифференцируем и объединяем дроби.}
\[
\begin{aligned}
D'(\omega)&=-\frac{2(R^2-2L/C)}{\omega^3}
-\frac{4}{\omega^5C^2}\\
&=\frac{2(2L-R^2C)C\omega^2-4}{C^2\omega^5}.
\end{aligned}
\]
При $\omega>0$ знаменатель положителен, поэтому определяющим остаётся числитель.

\textbf{Шаг 4. Проверяем существование положительного критического значения.}
Если
\[
\boxed{\ 2L>R^2C\ },
\]
то из $D'(\omega)=0$ следует
\[
\boxed{\omega_*=\sqrt{\frac{2}{C(2L-R^2C)}}}.
\]
До $\omega_*$ производная $D'$ отрицательна, после~--- положительна. Следовательно, $D$ минимальна, а $U_L$ максимальна при $\omega_*$. Если $2L\le R^2C$, положительного нуля производной нет: напряжение монотонно возрастает и стремится к $U$ при $\omega\to+\infty$, но конечная частота максимума отсутствует.

\textbf{Шаг 5. Подставляем числовые значения примера.} Пусть
\[
L=0{,}057\ \text{Гн},\quad C=2\cdot10^{-6}\ \text{Ф},\quad R=50\ \text{Ом}.
\]
Тогда $2L-R^2C=0{,}114-0{,}005=0{,}109>0$ и
\[
\omega_*=\sqrt{\frac{2}{2\cdot10^{-6}\cdot0{,}109}}
\approx\boxed{3028{,}9\ \text{рад/с}}.
\]
Если требуется обычная частота в герцах, применяем $\nu=\omega/(2\pi)$: $\nu_*\approx482{,}1$~Гц.

\section*{6. Как обнаружить ошибку до получения ответа}
\begin{enumerate}[label=\arabic*.]
\item \key{Проверьте степень при переносе под корень.} При делении всего знаменателя на $x^2$ выражение под корнем делится на $x^4$. Частая ошибка~--- сохранить $x^2$ вместо $x^4$.
\item \key{Проверьте раскрытие скобок.} В выражении $(a-b)^2$ средний член отрицателен. При сборе подобных слагаемых в задаче о волне получается коэффициент $(n-2)$, а не $(2-n)$.
\item \key{Проверьте производную степеней.} Например,
\[
\left(\frac{1}{\omega^4}\right)'=-\frac{4}{\omega^5},
\qquad
\left(\frac{R^2}{\omega^2}\right)'=-\frac{2R^2}{\omega^3}.
\]
\item \key{Проверяйте результат на невозможность.} Уравнение $\omega^4=-a$ при $a>0$ не имеет действительных решений. Такой результат может означать ошибку в знаке, неверную алгебру или отсутствие допустимого критического значения. Нужно вернуться к исходной формуле и перепроверить переходы.
\item \key{Не теряйте смысл параметра.} Условие $0<n<0{,}5$ задаёт диапазон, а не конкретное значение. Для вычисления одного числа требуется указать $n$.
\item \key{Отличайте максимум от предела.} Если функция строго возрастает на $x>0$ и только стремится к верхней границе при $x\to+\infty$, наибольшего значения на этой области у неё нет.
\item \key{Проверяйте физический смысл.} Период $\tau$, циклическая частота $\omega$ и электрические параметры в рассмотренных моделях положительны. Отрицательный корень времени или частоты исключается с объяснением.
\end{enumerate}

\section*{7. Финальная самопроверка решения}
\checkit{Я выписал исходную функцию без изменений и указал, что именно требуется максимизировать.}
\checkit{При преобразовании дроби сохранил все ограничения: $x>0$, ненулевой знаменатель и положительное подкоренное выражение.}
\checkit{Записал вспомогательную функцию $D(x)$ и объяснил, почему ищу именно её минимум.}
\checkit{Нашёл производную, приравнял числитель к нулю только после проверки знака знаменателя.}
\checkit{Исследовал знак производной по обе стороны критических точек и учёл границы области.}
\checkit{Подставил заданные параметры, проверил единицы и записал итоговый ответ с пояснением.}

\clearpage
\thispagestyle{fancy}
\begin{center}
{\Large\textcolor{darkaccent}{Алгоритм: поиск максимума сложной дроби}}\par
\vspace{5mm}
\end{center}
\begin{center}
\begin{tikzpicture}[
 node distance=7.5mm,
 process/.style={rectangle,rounded corners=2mm,draw=accent,line width=.8pt,align=center,text width=98mm,inner sep=6pt,font=\small},
 terminus/.style={ellipse,draw=darkaccent,line width=.9pt,align=center,text width=88mm,inner sep=5pt,font=\small},
 decision/.style={diamond,aspect=4.5,draw=accent,line width=.8pt,align=center,text width=50mm,inner sep=3pt,font=\small},
 arr/.style={-{Stealth[length=2.7mm]},draw=darkaccent,line width=.8pt}
]
\node[terminus] (start) {Дана функция $M(x)=K/\sqrt{D(x)}$, $K>0$};
\node[process,below=of start] (domain) {Определите область $x$; убедитесь, что $D(x)>0$};
\node[process,below=of domain] (simp) {Упростите $D(x)$: раскройте скобки, сократите степени и дроби};
\node[process,below=of simp] (deriv) {Найдите $D'(x)$, критические точки и знаки производной};
\node[decision,below=12mm of deriv] (min) {Достигает ли $D$ минимума?};
\node[process,below left=12mm and -9mm of min,text width=47mm] (yes) {Да: в точке минимума $D$ достигается максимум $M$};
\node[process,below right=12mm and -9mm of min,text width=47mm] (no) {Нет: изучите границы и пределы; максимума может не быть};
\node[terminus] (end) at ($(yes.south)!0.5!(no.south)+(0,-26mm)$) {Проверьте ограничения, параметры и единицы измерения};
\draw[arr] (start)--(domain);
\draw[arr] (domain)--(simp);
\draw[arr] (simp)--(deriv);
\draw[arr] (deriv)--(min);
\draw[arr] (min.south west) -- node[above left,font=\scriptsize]{да} (yes.north);
\draw[arr] (min.south east) -- node[above right,font=\scriptsize]{нет} (no.north);
\draw[arr] (yes.south) |- (end.west);
\draw[arr] (no.south) |- (end.east);
\end{tikzpicture}
\end{center}

\clearpage
\section*{Тренировочный блок}
\textcolor{darkaccent}{\textbf{Путь А. Десять задач от базовых преобразований к полноценной оптимизации.}}

Выполняйте задания последовательно. В каждой задаче приводите необходимые преобразования, записывайте ограничения, а при поиске максимума или минимума обосновывайте характер экстремума. В тренировочном блоке даны только условия; итоговые ответы размещены отдельно.

\subsection*{Базовый уровень}
\task{1}{Преобразование дроби с переменной под корнем}
При описании амплитуды некоторого колебательного процесса получена функция
\[
H(x)=\frac{18x^2}{\sqrt{(25-x^2)^2+6x^2}},\qquad x>0.
\]
Сначала укажите, почему при выполнении преобразований допустимо делить числитель и знаменатель на $x^2$. Затем представьте $H(x)$ в виде $18/\sqrt{F(x)}$, где числитель новой дроби не зависит от $x$. Раскройте квадрат разности, приведите подобные слагаемые и запишите $F(x)$ в виде суммы степеней $x$ с отрицательными показателями и постоянного числа. Проверьте результат обратным умножением на $x^4$ под корнем.

\task{2}{Раскрытие сложного квадрата и контроль показателей}
Вспомогательная функция задачи об электрической цепи имеет вид
\[
F(\omega)=\frac{R^2}{\omega^2}
+\frac{1}{\omega^2}\left(\omega L-\frac{1}{\omega C}\right)^2,
\qquad \omega>0.
\]
Для отдельного тренировочного расчёта примите $R=6$, $L=3$, $C=\tfrac12$ в согласованных условных единицах. Раскройте квадрат разности, умножьте каждый полученный член на $1/\omega^2$ и приведите подобные слагаемые. Запишите итоговую функцию только через степени $\omega$, найдите её производную и установите, положительна она или отрицательна на области $\omega>0$. Отдельно проверьте, где в результате находится четвёртая степень $\omega$.

\task{3}{Производная и проверка минимума}
Для вспомогательной функции
\[
F(t)=\frac{64}{t^4}-\frac{12}{t^2}+1,\qquad t>0,
\]
найдите производную и приведите её к одной дроби. Определите все допустимые значения $t$, при которых производная обращается в нуль. Исследуйте её знак слева и справа от найденной точки, определите промежутки возрастания и убывания функции, а затем вычислите наименьшее значение $F(t)$. Объясните, почему найденная точка действительно даёт минимум на всей области $t>0$.

\subsection*{Средний уровень}
\task{4}{Максимум дроби без прямого дифференцирования}
В некоторой модели эффективность устройства зависит от регулируемого положительного параметра $x$ по закону
\[
P(x)=\frac{120}{\sqrt{2+(x-4)^2}},\qquad x>0.
\]
Найдите значение параметра, при котором эффективность наибольшая, и вычислите это наибольшее значение. Обоснуйте ответ через свойства квадрата и зависимость дроби от знаменателя. Прямое дифференцирование исходной функции не требуется. Убедитесь, что выбранное значение допустимо.

\task{5}{Волновая задача с известным коэффициентом}
Лайнер совершает вынужденные поперечные колебания. Амплитуда описывается формулой
\[
A(\tau)=\frac{K\tau^2}{\sqrt{(T^2-\tau^2)^2+nT^2\tau^2}},\qquad \tau>0.
\]
В рамках данной модели $K=5$, период собственных колебаний равен $T=6$~с, а коэффициент $n=\tfrac14$. Найдите период приходящей волны, при котором амплитуда будет наибольшей, и само наибольшее значение амплитуды. Для обоснования преобразуйте функцию в обратную квадратному корню, найдите и исследуйте производную подкоренного выражения. Отрицательные значения периода исключите с указанием причины.

\task{6}{Оптимальный период при неизвестном параметре}
Амплитуда колебаний лайнера задаётся той же зависимостью
\[
A(\tau)=\frac{K\tau^2}{\sqrt{(T^2-\tau^2)^2+nT^2\tau^2}},\qquad \tau>0,
\]
где $K>0$ постоянно, $T=14$~с и известно только, что $0<n<0{,}5$. Найдите формулу оптимального периода $\tau_*$ через $n$. Затем определите точный открытый интервал, в котором может находиться $\tau_*$ при всех допустимых $n$. Объясните, можно ли назвать одно конкретное число секунд без дополнительной информации о $n$, и почему границы интервала не следует включать в ответ.

\task{7}{Период и амплитуда с числовой проверкой}
Для модели вынужденных колебаний задано
\[
A(\tau)=\frac{12\tau^2}{\sqrt{(100-\tau^2)^2+20\tau^2}},\qquad \tau>0.
\]
Период измеряется в секундах, амплитуда~--- в градусах. Установите, каким значениям $T$ и $n$ из общей формулы соответствует заданная зависимость. Найдите период наибольшей амплитуды и её величину. Выведите точные ответы и приближённые значения с точностью до сотых. Проверьте, что полученный период больше периода собственных колебаний.

\subsection*{Повышенный уровень}
\task{8}{Резонанс в цепи переменного тока}
Для катушки в последовательной электрической цепи справедлива зависимость
\[
U_L(\omega)=\frac{UL\omega}{\sqrt{R^2+(\omega L-1/(\omega C))^2}},
\qquad \omega>0.
\]
Амплитуда напряжения источника $U=200$~В, индуктивность катушки $L=0{,}05$~Гн, ёмкость $C=10^{-5}$~Ф, сопротивление $R=20$~Ом. Найдите положительную циклическую частоту, при которой напряжение на катушке максимально, а также само максимальное напряжение. Докажите наличие максимума с помощью производной упрощённого подкоренного выражения. Частоту округлите до десятых радиан в секунду, напряжение~--- до десятых вольта.

\task{9}{Случай, когда конечного максимума нет}
При анализе последовательной электрической цепи получена формула
\[
U_L(\omega)=\frac{UL\omega}{\sqrt{R^2+(\omega L-1/(\omega C))^2}},\qquad \omega>0.
\]
Параметры цепи таковы: $U=100$~В, $L=0{,}02$~Гн, $C=2\cdot10^{-5}$~Ф, $R=50$~Ом. Проверьте условие существования положительной стационарной частоты. Затем исследуйте знак производной вспомогательного выражения для всех $\omega>0$ и опишите поведение $U_L(\omega)$. Определите, существует ли конечная частота достижения максимума, и найдите предельное напряжение при неограниченном увеличении частоты.

\task{10}{Полная оптимизация с переходом к герцам}
Источник переменного тока подключён к последовательной цепи с сопротивлением $R$, катушкой индуктивности $L$ и конденсатором ёмкости $C$. Амплитуда напряжения на катушке равна
\[
U_L(\omega)=\frac{UL\omega}{\sqrt{R^2+(\omega L-1/(\omega C))^2}},\qquad \omega>0.
\]
Известно, что $U=311$~В, $L=0{,}057$~Гн, $C=2\cdot10^{-6}$~Ф, $R=50$~Ом. Получите формулу для оптимальной циклической частоты, подставьте данные и определите её числовое значение. Переведите результат в обычную частоту $\nu=\omega/(2\pi)$ в герцах. Дополнительно найдите амплитуду напряжения на катушке в точке максимума. Обоснуйте, почему найденный экстремум является максимумом и почему величина напряжения на катушке может превосходить напряжение источника.

\clearpage
\section*{Ответы к тренировочному блоку}
\small
\begin{longtable}{@{}p{11mm}p{\dimexpr\linewidth-18mm\relax}@{}}
\toprule
\textbf{№} & \textbf{Ответ для самопроверки}\\
\midrule
\endhead
1 & $F(x)=\dfrac{625}{x^4}-\dfrac{44}{x^2}+1$; $H(x)=18/\sqrt{F(x)}$, $x>0$.\\[5pt]
2 & $F(\omega)=9+\dfrac{24}{\omega^2}+\dfrac4{\omega^4}$; $F'(\omega)=-\dfrac{48}{\omega^3}-\dfrac{16}{\omega^5}<0$ при $\omega>0$.\\[5pt]
3 & $F'(t)=\dfrac{24t^2-256}{t^5}$. Минимум при $t=\sqrt{\dfrac{32}{3}}=4\sqrt{\dfrac23}$; $F_{\min}=\dfrac7{16}$. До этой точки $F$ убывает, после возрастает.\\[5pt]
4 & $x=4$; $P_{\max}=60\sqrt2\approx84{,}85$.\\[5pt]
5 & $\tau_*=6\sqrt{\dfrac87}\approx6{,}41$~с; $A_{\max}=\dfrac{40}{\sqrt{15}}\approx10{,}33$.\\[5pt]
6 & $\tau_*=14\sqrt{\dfrac2{2-n}}$~с; $\displaystyle14<\tau_*<\frac{28}{\sqrt3}\ \text{с}$. Одного числа при неизвестном $n$ нет.\\[5pt]
7 & $T=10$~с, $n=\dfrac15$; $\tau_*=\dfrac{10\sqrt{10}}3\approx10{,}54$~с; $A_{\max}=\dfrac{120}{\sqrt{19}}\approx27{,}53^\circ$.\\[5pt]
8 & $\omega_*=\dfrac{2500}{\sqrt3}\approx1443{,}4$~рад/с; $U_{L,\max}=\dfrac{5000}{7}\approx714{,}3$~В.\\[5pt]
9 & $2L-R^2C=0{,}04-0{,}05<0$. Конечной частоты максимума нет; $U_L(\omega)$ строго возрастает и стремится к $100$~В снизу.\\[5pt]
10 & $\omega_*=\sqrt{\dfrac{2}{C(2L-R^2C)}}\approx3028{,}9$~рад/с; $\nu_*\approx482{,}1$~Гц; $U_{L,\max}\approx1061{,}8$~В.\\[2pt]
\bottomrule
\end{longtable}
\normalsize
\vspace{7mm}
\textbf{Памятка.} Отрицательные значения времени и циклической частоты в этих моделях исключаются. Перед подстановкой параметров проверьте, что знаменатели отличны от нуля, подкоренные выражения положительны, а стационарная точка действительно соответствует минимуму вспомогательной функции.

\vfill
\begin{center}\small\textcolor{muted}{Лёвин Артём Александрович эксклюзивно для Матвея Горбачева}\end{center}
\end{document}
